\documentclass[11pt]{article}
\usepackage{amsmath,amsthm,amssymb}
\usepackage[margin=1in]{geometry}

\newtheorem{theorem}{Theorem}
\newtheorem{lemma}{Lemma}
\newtheorem{proposition}{Proposition}
\newtheorem{corollary}{Corollary}
\theoremstyle{definition}
\newtheorem{remark}{Remark}
\newtheorem{definition}{Definition}

\DeclareMathOperator{\Des}{Des}
\DeclareMathOperator{\maj}{maj}
\DeclareMathOperator{\comaj}{comaj}
\DeclareMathOperator{\SYT}{SYT}
\DeclareMathOperator{\SDT}{SDT}
\DeclareMathOperator{\ord}{ord}
\DeclareMathOperator{\core}{core}
\DeclareMathOperator{\hght}{ht}
\newcommand{\ZZ}{\mathbb{Z}}
\newcommand{\la}{\lambda}
\newcommand{\inner}[1]{\left\langle #1 \right\rangle}

\title{The leading coefficient of the graded $2$-core law:\\
a refined $2$-quotient character identity and the closed form
$S(n,c)=\varepsilon(c_0)\,2^{-c}\,e_c(W_n)$}
\author{Clio}
\date{2026-06-02}

\begin{document}
\maketitle

\begin{abstract}
For $\la\vdash n$ let $G_\la(q)=\sum_{T\in\SYT(\la)}q^{s(T)}$ be the graded
branch-exponent enumerator, $s(T)=\sum_{i\in\Des(T)}w_i$ with parity-twisted
weights $w_i=(2i-1)[\,n-i\text{ odd}\,]$. The graded $2$-core law
($\ord_{q=-1}G_\la=c:=\lfloor c_0/2\rfloor$, $c_0=|2\text{-}\core(\la)|$) is
proved, and its leading Taylor coefficient at $q=-1$ was reduced to one
$\la$-dependent character value $\chi^\la(2^{w}1^{c_0})$, $w=(n-c_0)/2$. Here I
close that last piece. I prove the \emph{refined $2$-quotient character
identity}
\[
\chi^\la\!\big(2^{w}1^{c_0}\big)
   =\varepsilon(c_0)\,(-1)^{n(\la)}\,f^{\la^{(0)}}f^{\la^{(1)}}
     \binom{w}{|\la^{(0)}|}\,f^{\,2\text{-}\core(\la)},
\]
where $(\la^{(0)},\la^{(1)})$ is the $2$-quotient and the sign
$\varepsilon(c_0)=(-1)^{n(\delta_k)}=(-1)^{\binom{k+1}{3}}$ depends only on the
$2$-core $\delta_k=(k,k-1,\dots,1)$ (so on $c_0=\binom{k+1}{2}$), being $-1$
exactly when $k\equiv 2\pmod 4$. The proof is one iterated application of the
Murnaghan--Nakayama rule by dominoes: domino removal preserves the $2$-core, so
the strip-tableau sum collapses onto the single core term, whose signed count
factors through the Stanton--White correspondence with a sign that is constant by
a row-parity invariant. Assembling with the master identity yields the leading
coefficient in closed form,
\[
a_c=(-1)^{n(\la)}\varepsilon(c_0)\,2^{-c}\,e_c(W_n)\,
     f^{\la^{(0)}}f^{\la^{(1)}}\binom{w}{|\la^{(0)}|}\,f^{\,2\text{-}\core(\la)},
\qquad
S(n,c)=\varepsilon(c_0)\,2^{-c}\,e_c(W_n),
\]
with $W_n$ the arithmetic progression $\{4j+a:0\le j<N\}$
($(a,N)=(1,n/2)$ for $n$ even, $(3,(n-1)/2)$ for $n$ odd); I give $e_c(W_n)$ in
closed (Pochhammer / generalized-Stirling) form. The single-valuedness of
$S(n,c)$ across all $\la\vdash n$ with fixed $(n,c)$ is exactly the content of
the character identity. All formulas are verified by exact computation for every
partition of every $n\le 12$ ($270$ shapes, $0$ mismatches).
\end{abstract}

\section{Statement and setup}

I keep the notation of the graded $2$-core law (``06-01''). Throughout
$\la\vdash n$; $\SYT(\la)$, $\Des(T)$, $\comaj(T)=\sum_{i\in\Des(T)}(n-i)$,
$n(\la)=\sum_i(i-1)\la_i$, $f^\mu=\#\SYT(\mu)$ are standard. The branch weights
are $w_i=(2i-1)\,[\,n-i\text{ odd}\,]$, the \emph{active set} is
$A=\{k\in[n-1]:n-k\text{ odd}\}$ with $|A|=m=\lfloor n/2\rfloor$, and
\[
G_\la(q)=\sum_{T\in\SYT(\la)}q^{\,s(T)},\qquad s(T)=\sum_{i\in\Des(T)}w_i .
\]
Every $\la$ has a $2$-core, which is a staircase
$\delta_k=(k,k-1,\dots,1)$ of size $c_0:=|2\text{-}\core(\la)|=\binom{k+1}{2}$,
and a $2$-quotient $(\la^{(0)},\la^{(1)})$ with
$|\la^{(0)}|+|\la^{(1)}|=w:=(n-c_0)/2$ (the number of dominoes). Set
$c=\lfloor c_0/2\rfloor$, so $c_0=2c+e$ with $e=n\bmod 2$, and note
$m-c=w$.

\paragraph{What is already proved (06-01).}
The order law $\ord_{q=-1}G_\la=c$ holds, and the first nonzero Taylor
coefficient $a_c$ in the expansion $G_\la(q)=\sum_k a_k(q+1)^k$ is
\begin{equation}\label{eq:ac-master}
  a_c=(-1)^{c}\,e_c(W_n)\,\Big(-\tfrac12\Big)^{c}\,\chi^\la\!\big(2^{\,w}1^{\,c_0}\big)
     =2^{-c}\,e_c(W_n)\,\chi^\la\!\big(2^{\,w}1^{\,c_0}\big),
\end{equation}
where
\begin{equation}\label{eq:Wn}
  W_n=\{\,2k-1:k\in A\,\}
\end{equation}
and $e_c$ is the elementary symmetric polynomial. (The two forms agree because
$(-1)^c(-\tfrac12)^c=2^{-c}$.) Equation~\eqref{eq:ac-master} is line-for-line the
output of the 06-01 master identity; the only object in it that still depended on
$\la$ in an unresolved way is the character value $\chi^\la(2^{w}1^{c_0})$ --- a
class with $w$ transpositions and $c_0$ fixed points.

\paragraph{Goal.}
Evaluate $\chi^\la(2^{w}1^{c_0})$ in closed form (Theorem~\ref{thm:char}),
deduce the leading coefficient and $S(n,c)$ (Theorem~\ref{thm:Snc}), pin the sign
$\varepsilon(c_0)$ (Proposition~\ref{prop:sign}), and give $e_c(W_n)$ in closed
form (Proposition~\ref{prop:ec}).

\section{The refined $2$-quotient character identity}\label{sec:char}

\begin{theorem}\label{thm:char}
Let $\la\vdash n$ have $2$-core $\delta_k$ of size $c_0=\binom{k+1}{2}$,
$2$-quotient $(\la^{(0)},\la^{(1)})$, and $w=(n-c_0)/2$ dominoes. Then
\[
  \chi^\la\!\big(2^{\,w}1^{\,c_0}\big)
   =\varepsilon(c_0)\,(-1)^{n(\la)}\,
     f^{\la^{(0)}}f^{\la^{(1)}}\binom{w}{|\la^{(0)}|}\,f^{\delta_k},
   \qquad
   \varepsilon(c_0)=(-1)^{n(\delta_k)}\in\{\pm1\}.
\]
In particular the value is nonzero, and depends on $\la$ only through
$(-1)^{n(\la)}$ and the $2$-quotient data; the sign $\varepsilon(c_0)$ depends
only on $c_0$.
\end{theorem}

This generalizes the $c_0\in\{0,1\}$ case (the reversal/longest-element class
$w_0$), proved in ``06-03'', where $\delta_k\in\{\varnothing,(1)\}$,
$f^{\delta_k}=1$, and $\varepsilon=+1$. The proof is the same engine --- iterated
domino Murnaghan--Nakayama --- carried out over a nonempty staircase core.

\subsection{Collapse onto the core: only $\mu=\delta_k$ survives}

\begin{lemma}\label{lem:collapse}
Apply the Murnaghan--Nakayama (MN) rule to $\chi^\la$ on the class
$2^{w}1^{c_0}$, removing the $w$ parts of size $2$ first and the $c_0$ fixed
points last. Then
\[
  \chi^\la\!\big(2^{w}1^{c_0}\big)=N(\la)\cdot f^{\delta_k},
  \qquad
  N(\la)=\sum_{\substack{\la=\nu_0\supset\nu_1\supset\cdots\supset\nu_w=\delta_k\\
                          \nu_{i-1}/\nu_i\ \mathrm{a\ domino}}}
         \ \prod_{i=1}^{w}(-1)^{\hght(\nu_{i-1}/\nu_i)-1},
\]
the signed count of domino-strip removal sequences from $\la$ down to its
$2$-core, where $\hght$ is the number of rows a strip occupies.
\end{lemma}

\begin{proof}
The MN rule, with the cycle parts ordered as $w$ twos then $c_0$ ones, gives
\[
  \chi^\la\!\big(2^{w}1^{c_0}\big)
  =\sum_{\la=\nu_0\supset\cdots\supset\nu_w}
    \Big(\prod_{i=1}^w(-1)^{\hght(\nu_{i-1}/\nu_i)-1}\Big)\,
    \chi^{\nu_w}\!\big(1^{c_0}\big),
\]
the sum over chains in which each skew $\nu_{i-1}/\nu_i$ is a border strip
(rim hook) of size $2$, i.e.\ a domino, so $|\nu_w|=n-2w=c_0$; and
$\chi^{\nu_w}(1^{c_0})=f^{\nu_w}$ (removing $c_0$ single cells, all of height
$1$, contributes only $+1$'s, and counts SYT). It remains to see that the only
$\nu_w\vdash c_0$ that occurs is $\delta_k$.

Removing a domino (a $2$-rim-hook) does not change the $2$-core of a partition
(James--Kerber, \emph{The Representation Theory of the Symmetric Group}, \S2.7).
Hence every $\nu_i$ in such a chain has the same $2$-core as $\la$, namely
$\delta_k$; in particular $2\text{-}\core(\nu_w)=\delta_k$, so
$|2\text{-}\core(\nu_w)|=c_0=|\nu_w|$. A partition whose $2$-core has its own
size is equal to its $2$-core (the core is reached by removing dominoes, each
lowering the size by $2$; equal size forces zero removals). Therefore
$\nu_w=\delta_k$, and $\chi^{\nu_w}(1^{c_0})=f^{\delta_k}$ in every surviving
term. Factoring it out gives the claim.
\end{proof}

A removal chain $\la=\nu_0\supset\cdots\supset\nu_w=\delta_k$ is the same datum
as a \emph{standard domino tableau} of the skew shape $\la/\delta_k$: tile
$\la/\delta_k$ by $w$ dominoes and number them $1,\dots,w$ by reverse removal
order. The MN weight of a chain is
\[
  \prod_{i=1}^w(-1)^{\hght(\nu_{i-1}/\nu_i)-1}=(-1)^{v(D)},
\]
where $v(D)$ is the number of \emph{vertical} dominoes of the underlying tiling
$D$ (a horizontal domino has $\hght=1$, a vertical one has $\hght=2$). This weight
depends only on the tiling $D$, not on the numbering. So
\begin{equation}\label{eq:Nsum}
  N(\la)=\sum_{D}\;(-1)^{v(D)}\,\#\{\text{standard fillings of }D\},
\end{equation}
the sum over domino tilings $D$ of $\la/\delta_k$.

\subsection{The sign is constant: a row-parity invariant}

\begin{lemma}\label{lem:sign}
For every domino tiling $D$ of the skew shape $\la/\delta_k$,
\[
  v(D)\equiv n(\la)+n(\delta_k)\pmod 2 .
\]
In particular $(-1)^{v(D)}$ is the same for all tilings, equal to
$(-1)^{n(\la)+n(\delta_k)}$.
\end{lemma}

\begin{proof}
Colour a cell $(i,j)$ by the parity of its row index $i$. A horizontal domino
occupies $(i,j),(i,j+1)$ --- two cells of the \emph{same} row-parity; a vertical
domino occupies $(i,j),(i+1,j)$ --- one cell of each row-parity. Let $a$ be the
number of cells of $\la/\delta_k$ lying in odd rows. Counting odd-row cells by
domino type,
\[
  a=v(D)\cdot 1+h_{\mathrm{odd}}(D)\cdot 2,
\]
where $h_{\mathrm{odd}}(D)$ is the number of horizontal dominoes in odd rows
(each vertical domino contributes exactly one odd-row cell; a horizontal domino
contributes $2$ if it sits in an odd row and $0$ otherwise). Hence
$v(D)\equiv a\pmod 2$, and the right-hand side $a$ depends only on
$\la/\delta_k$, not on $D$.

It remains to evaluate $a\bmod 2$. For any partition $\mu$, writing
$b_\mu=\sum_{i\text{ even}}\mu_i$ for the number of cells of $\mu$ in even rows,
we have
\[
  n(\mu)=\sum_i(i-1)\mu_i\equiv\sum_{i\text{ even}}\mu_i=b_\mu\pmod 2,
\]
since $i-1$ is odd exactly when $i$ is even. The number of odd-row cells of $\mu$
is $|\mu|-b_\mu$. Therefore
\[
  a=\big(n-b_\la\big)-\big(c_0-b_{\delta_k}\big)
   =(n-c_0)-b_\la+b_{\delta_k}
   =2w-b_\la+b_{\delta_k}
   \equiv b_\la+b_{\delta_k}
   \equiv n(\la)+n(\delta_k)\pmod 2,
\]
which gives the claim.
\end{proof}

\begin{remark}
Lemma~\ref{lem:sign} is the skew-shape form of the spin/$2$-quotient parity
invariant; over an empty core it recovers the ``06-03'' statement
$v(D)\equiv n(\la)\pmod2$. The row-parity count makes both the constancy and the
value visible at once, with no appeal to the Stanton--White bijection.
\end{remark}

\subsection{The count: Stanton--White}

\begin{lemma}\label{lem:count}
The number of standard domino tableaux of $\la$ (equivalently, of standard
fillings of all domino tilings of $\la/\delta_k$) is
\[
  \#\SDT(\la)=f^{\la^{(0)}}\,f^{\la^{(1)}}\,\binom{w}{|\la^{(0)}|}.
\]
\end{lemma}

\begin{proof}
This is the Stanton--White correspondence (James--Kerber \S2.7; Stanton--White,
\emph{J.\ Combin.\ Theory A} 40 (1985)). It is a bijection between standard
domino tableaux of $\la$ and triples $(P^{(0)},P^{(1)},\sigma)$, where
$P^{(r)}\in\SYT(\la^{(r)})$ and $\sigma$ is a shuffle distributing the $w$ time
labels $\{1,\dots,w\}$ among the two quotient components, $|\la^{(0)}|$ to the
first and $|\la^{(1)}|$ to the second. The number of such shuffles is
$\binom{w}{|\la^{(0)}|}$, giving the product. (The bijection is purely
combinatorial and holds for every $\la$, irrespective of its core.)
\end{proof}

\subsection{Assembly of Theorem~\ref{thm:char}}

By \eqref{eq:Nsum} and Lemma~\ref{lem:sign}, the sign is constant over tilings,
so
\[
  N(\la)=(-1)^{n(\la)+n(\delta_k)}\sum_D\#\{\text{standard fillings of }D\}
        =(-1)^{n(\la)+n(\delta_k)}\,\#\SDT(\la).
\]
Combining with Lemma~\ref{lem:collapse} and Lemma~\ref{lem:count},
\[
  \chi^\la\!\big(2^{w}1^{c_0}\big)
   =N(\la)\,f^{\delta_k}
   =(-1)^{n(\delta_k)}(-1)^{n(\la)}\,
     f^{\la^{(0)}}f^{\la^{(1)}}\binom{w}{|\la^{(0)}|}\,f^{\delta_k}.
\]
Setting $\varepsilon(c_0)=(-1)^{n(\delta_k)}$ --- which depends only on $k$,
hence only on $c_0=\binom{k+1}{2}$ --- proves Theorem~\ref{thm:char}. The value
is nonzero because $f^{\la^{(0)}},f^{\la^{(1)}},f^{\delta_k}\ge1$ and
$\binom{w}{|\la^{(0)}|}\ge1$. \qquad$\square$

\section{The leading coefficient and $S(n,c)$}\label{sec:Snc}

\begin{theorem}\label{thm:Snc}
With $W_n$ as in \eqref{eq:Wn}, the leading Taylor coefficient of $G_\la$ at
$q=-1$ is
\[
  a_c=(-1)^{n(\la)}\,\varepsilon(c_0)\,2^{-c}\,e_c(W_n)\,
       f^{\la^{(0)}}f^{\la^{(1)}}\binom{w}{|\la^{(0)}|}\,f^{\delta_k}.
\]
Equivalently, writing
$a_c=(-1)^{n(\la)}f^{\la^{(0)}}f^{\la^{(1)}}\binom{w}{|\la^{(0)}|}f^{\delta_k}\,S(n,c)$,
\[
  \boxed{\,S(n,c)=\varepsilon(c_0)\,2^{-c}\,e_c(W_n)\,}.
\]
The scalar $S(n,c)$ depends only on $(n,c)$; this single-valuedness over all
$\la\vdash n$ with $2$-core size $c_0=2c+e$ is precisely Theorem~\ref{thm:char}.
\end{theorem}

\begin{proof}
Substitute Theorem~\ref{thm:char} into the master form \eqref{eq:ac-master}:
\[
  a_c=2^{-c}e_c(W_n)\,
      \varepsilon(c_0)(-1)^{n(\la)}f^{\la^{(0)}}f^{\la^{(1)}}
      \binom{w}{|\la^{(0)}|}f^{\delta_k}.
\]
Dividing by
$(-1)^{n(\la)}f^{\la^{(0)}}f^{\la^{(1)}}\binom{w}{|\la^{(0)}|}f^{\delta_k}$ (each
factor nonzero) isolates $S(n,c)=\varepsilon(c_0)2^{-c}e_c(W_n)$, manifestly a
function of $(n,c)$ alone.
\end{proof}

\begin{remark}[Two conventions for $S$]
Some notes fold $f^{\delta_k}$ into $S$, writing
$a_c=(-1)^{n(\la)}f^{\la^{(0)}}f^{\la^{(1)}}\binom{w}{|\la^{(0)}|}\,\widetilde
S(n,c)$ with $\widetilde S(n,c)=\varepsilon(c_0)2^{-c}e_c(W_n)f^{\delta_k}$; then
$\widetilde S(n,1)=-\binom n2$ (since $f^{\delta_2}=f^{(2,1)}=2$). The convention
of Theorem~\ref{thm:Snc} keeps $f^{\delta_k}$ separate, giving the cleaner
$S(n,1)=-\tfrac12\binom n2$. Both are equivalent statements about the same
$a_c$.
\end{remark}

\section{The sign $\varepsilon(c_0)$ in closed form}\label{sec:sign}

\begin{proposition}\label{prop:sign}
For the staircase core $\delta_k=(k,k-1,\dots,1)$ with $c_0=\binom{k+1}{2}$,
\[
  \varepsilon(c_0)=(-1)^{n(\delta_k)}=(-1)^{\binom{k+1}{3}},
\]
and $\varepsilon(c_0)=-1$ if and only if $k\equiv 2\pmod 4$.
\end{proposition}

\begin{proof}
Compute $n(\delta_k)$. With $\delta_{k,i}=k-i+1$,
\[
  n(\delta_k)=\sum_{i=1}^{k}(i-1)(k-i+1)
            =\sum_{j=0}^{k-1}j\,(k-j)
            =k\frac{(k-1)k}{2}-\frac{(k-1)k(2k-1)}{6}
            =\frac{(k-1)k(k+1)}{6}=\binom{k+1}{3}.
\]
Hence $\varepsilon(c_0)=(-1)^{\binom{k+1}{3}}$. For the parity criterion, by
Lucas' theorem $\binom{k+1}{3}$ is odd iff the binary digits of $3=(11)_2$ lie
under those of $k+1$, i.e.\ iff the last two bits of $k+1$ are both $1$, i.e.\
$k+1\equiv 3\pmod 4$, i.e.\ $k\equiv 2\pmod 4$.
\end{proof}

Thus $\varepsilon=+1$ for $k\in\{0,1,3,4,5,7,8,9,\dots\}$ and $\varepsilon=-1$
for $k\in\{2,6,10,\dots\}$; in $c$-terms this is $\varepsilon(c_0)=-1$ exactly at
$c\in\{1,10,\dots\}$ (i.e.\ $c_0\in\{3,21,\dots\}$), matching the observed
$\varepsilon=-1$ at $c=1$ and $\varepsilon=+1$ at $c\in\{0,3,5,7\}$.

\section{Closed form for $e_c(W_n)$}\label{sec:ec}

\begin{proposition}\label{prop:ec}
The set $W_n=\{2k-1:k\in A\}$ is the arithmetic progression
\[
  W_n=\{\,4j+a:0\le j<N\,\},\qquad
  (a,N)=\begin{cases}(1,\ n/2),&n\text{ even},\\[2pt](3,\ (n-1)/2),&n\text{ odd}.\end{cases}
\]
Consequently its elementary symmetric polynomials have generating function
\[
  \sum_{c\ge0}e_c(W_n)\,t^{c}
   =\prod_{j=0}^{N-1}\big(1+(4j+a)t\big)
   =(4t)^{N}\Big(\tfrac{a}{4}+\tfrac{1}{4t}\Big)_{N},
\]
where $(z)_N=z(z+1)\cdots(z+N-1)$ is the Pochhammer (rising-factorial) symbol;
equivalently, via the unsigned Stirling numbers of the first kind
$\big[{N\atop \kappa}\big]$,
\[
  e_c(W_n)=\sum_{\kappa=0}^{N}\Big[{N\atop \kappa}\Big]\,
            4^{\,N-\kappa}\,a^{\,c-N+\kappa}\,\binom{\kappa}{\,c-N+\kappa\,},
\]
the sum running over $\kappa$ with $0\le c-N+\kappa\le \kappa$.
\end{proposition}

\begin{proof}
If $n$ is even then $n-k$ is odd iff $k$ is odd, so $A=\{1,3,\dots,n-1\}$ and
$W_n=\{1,5,\dots,2n-3\}=\{4j+1:0\le j<n/2\}$. If $n$ is odd then $n-k$ is odd iff
$k$ is even, so $A=\{2,4,\dots,n-1\}$ and
$W_n=\{3,7,\dots,2n-3\}=\{4j+3:0\le j<(n-1)/2\}$. This is the displayed
progression with common difference $4$.

The first factorization is the definition of $e_c$ as the coefficients of
$\prod_{w\in W_n}(1+wt)$. For the Pochhammer form, write
$1+(4j+a)t=4t\big(j+\tfrac{1+at}{4t}\big)$ and use
$\prod_{j=0}^{N-1}(j+z)=(z)_N$ with $z=\tfrac{1+at}{4t}=\tfrac{a}{4}+\tfrac1{4t}$.
For the Stirling form, expand $(z)_N=\sum_\kappa\big[{N\atop\kappa}\big]z^\kappa$
and substitute $z=\tfrac{a}{4}+\tfrac1{4t}$, i.e.
\[
  \prod_{j=0}^{N-1}\big(1+(4j+a)t\big)
   =\sum_{\kappa=0}^N\Big[{N\atop\kappa}\Big](4t)^{N-\kappa}(1+at)^{\kappa};
\]
reading off $[t^c]$ (the term needs $t^{N-\kappa}\cdot t^{\,c-N+\kappa}$ from
$(1+at)^\kappa$) gives the stated convolution.
\end{proof}

For instance $e_1(W_n)=\sum_{w\in W_n}w=\binom n2$ (whence
$S(n,1)=-\tfrac12\binom n2$), and the smallest higher values are
$e_3(W_6)=1\cdot5\cdot9=45$ ($S(6,3)=45/8$),
$e_5(W_{10})=1\cdot5\cdot9\cdot13\cdot17=9945$ ($S(10,5)=9945/32$).

\section{Computational verification}

All claims were checked by independent exact (rational) computation; the SYT side
(direct enumeration of $\SYT(\la)$ with $s(T)$, Taylor expansion of $G_\la$ about
$q=-1$) shares no code with the character/$2$-quotient side.

\begin{itemize}
\item \textbf{Leading coefficient (Theorem~\ref{thm:Snc}).} For all $\la\vdash n$,
$2\le n\le 12$ ($270$ shapes), the directly computed $a_c$ equals
$(-1)^{n(\la)}\varepsilon(c_0)2^{-c}e_c(W_n)f^{\la^{(0)}}f^{\la^{(1)}}
\binom{w}{|\la^{(0)}|}f^{\delta_k}$: \emph{$0$ mismatches}.
\item \textbf{Single-valuedness / character identity
(Theorem~\ref{thm:char}).} The reduced scalar
$a_c/\big((-1)^{n(\la)}f^{\la^{(0)}}f^{\la^{(1)}}\binom{w}{|\la^{(0)}|}
f^{\delta_k}\big)$ is constant across all $\la\vdash n$ of each fixed $(n,c)$,
$n\le12$, and equals $\varepsilon(c_0)2^{-c}e_c(W_n)$. Sample table
($S(n,c)$): $S(n,1)=-\tfrac12\binom n2$; $S(6,3)=45/8$, $S(8,3)=203/2$,
$S(10,3)=2565/4$, $S(12,3)=5115/2$; $S(10,5)=9945/32$, $S(12,5)=156057/16$.
\item \textbf{Sign (Proposition~\ref{prop:sign}).}
$\varepsilon(c_0)=(-1)^{n(\delta_k)}=(-1)^{\binom{k+1}{3}}$ for $k=0,\dots,6$:
$\varepsilon=+1$ except $k\in\{2,6\}$ where $\varepsilon=-1$ (matching
$c\in\{1,10\}$).
\item \textbf{$e_c$ closed forms (Proposition~\ref{prop:ec}).} The AP
description, the Pochhammer generating function, and the Stirling convolution all
reproduce $e_c(W_n)$ exactly for every $n\le 12$ and every $c$.
\end{itemize}

\paragraph{Scripts.}
\texttt{\~{}/projects/scratch/2026-06-05-Snc/} (notably \texttt{snc12.py} for the
single-valued $S(n,c)$ table, \texttt{verify\_closedform.py} for the full $a_c$
identity and the sign, \texttt{verify\_ec.py} for the $e_c$ closed forms), reusing
the $\SYT$/$2$-quotient/Murnaghan--Nakayama routines of
\texttt{\~{}/projects/scratch/prove-2026-06-01-graded-2core/}.

\section{Summary}

The graded $2$-core law $\ord_{q=-1}G_\la=c$ is now complete down to its leading
coefficient. The one remaining $\la$-dependent quantity, $\chi^\la(2^{w}1^{c_0})$,
is evaluated in closed form (Theorem~\ref{thm:char}) by iterated domino
Murnaghan--Nakayama: domino removal preserves the $2$-core, collapsing the
strip-tableau sum onto the single core term $\delta_k$; the resulting signed
domino-tableau count factors through Stanton--White, with a sign made constant by
a one-line row-parity invariant. Assembling with the master identity gives
\[
  a_c=(-1)^{n(\la)}\varepsilon(c_0)\,2^{-c}\,e_c(W_n)\,
       f^{\la^{(0)}}f^{\la^{(1)}}\binom{w}{|\la^{(0)}|}\,f^{\delta_k},
  \qquad
  S(n,c)=\varepsilon(c_0)\,2^{-c}\,e_c(W_n),
\]
with $\varepsilon(c_0)=(-1)^{\binom{k+1}{3}}$ ($-1$ iff $k\equiv2\bmod4$,
$c_0=\binom{k+1}{2}$) and $e_c(W_n)$ the elementary symmetric polynomial of the
arithmetic progression $\{4j+a:0\le j<N\}$, given in Pochhammer /
generalized-Stirling form. The single-valuedness of $S(n,c)$ \emph{is} the
character identity. No gaps remain.

\end{document}
